\clearpage
\begin{summarybox}
  \textbf{\textcolor{CSummary}{一句话记住}}

  \textbf{奇函数中，$x$ 和 $-x$ 的函数值互为相反数；而 $0=-0$，
  所以只要 $f(0)$ 存在，就只能等于 $0$。}

  \[
    \boxed{0\in D_f\ \text{且}\ f\text{为奇函数}\ \Longrightarrow\ f(0)=0}
  \]

  \medskip
  \textbf{\textcolor{CSummary}{两行证明}}
  \[
    f(-0)=-f(0)
    \ \Longrightarrow\ f(0)=-f(0)
    \ \Longrightarrow\boxed{f(0)=0}.
  \]

  \medskip
  \textbf{\textcolor{CSummary}{解题顺序}}
  \begin{enumerate}[leftmargin=*, itemsep=0.35em, topsep=0.3em]
    \item \textbf{看定义域：}$0$ 能否代入？
    \item \textbf{用已知性质：}若是奇函数，立即写 $f(0)=0$。
    \item \textbf{继续求参：}把 $x=0$ 代入表达式；需要时再验证 $f(-x)=-f(x)$。
  \end{enumerate}

  \medskip
  \textbf{\textcolor{CSummary}{三个不能忘}}
  \begin{itemize}[leftmargin=*, itemsep=0.4em, topsep=0.3em]
    \item $f(x)=\dfrac1x$ 是奇函数，但没有 $f(0)$。
    \item $f(x)=x^2$ 满足 $f(0)=0$，却不是奇函数。
    \item 奇函数若 $f(a)=0$，则 $f(-a)=0$；零点可以不止 $0$。
  \end{itemize}

  \medskip
  \textbf{\textcolor{CSummary}{一条更大的规律}}

  图像关于 $(a,b)$ 中心对称，且 $x=a$ 有定义时，
  \[
    \boxed{f(a)=b.}
  \]
  \textbf{本节结论正是对称中心取 $(0,0)$ 时的情况。}
\end{summarybox}
